\documentclass{article}
% Per un output più pulito della sola figura, decommenta la riga seguente
%\documentclass[border=10pt]{standalone}

% --- PACHETTI ---
\usepackage[T1]{fontenc}
\usepackage{amsmath}
\usepackage{amssymb}
\usepackage{graphicx}
\usepackage{caption}
\usepackage{subcaption}
\usepackage{pgfplots} % Carica anche TikZ e xcolor
\usepackage{fontspec}

% --- CONFIGURAZIONE ---
\pgfplotsset{compat=1.18}
\captionsetup{font=small, labelfont=bf, labelsep=period}
\setmainfont{Arial} % Richiede XeLaTeX o LuaLaTeX

\begin{document}

\begin{figure}[htbp]
\centering
\begin{tikzpicture}
\begin{axis}[
    width=12cm,
    height=8cm,
    xlabel={Latenza (ms)},
    ylabel={Sicurezza (ASSA Score)},
    xmin=10, xmax=100,
    ymin=0, ymax=100,
    colormap/viridis,
    colorbar,
    colorbar style={
        title={Throughput (TPS)},
        ylabel={TPS}
    },
    title={Trade-off Multi-Obiettivo nel Framework GIST\\{\small Frontiera di Pareto per Configurazioni Ottimali}}
]

% --- PLOT DI CONTORNO (CON GNUPLOT) ---
% RICORDA: Richiede Gnuplot installato e la compilazione con --shell-escape
\addplot3[
    contour gnuplot={
        levels={2000,3000,4000,5000,6000,7000,8000,9000},
        labels=false
    },
    thick
]
{10000*exp(-0.02*x)*(1+0.005*y)};

% --- PUNTI DELLE CONFIGURAZIONI ---
\addplot[only marks, mark=*, mark size=3pt, black] coordinates {
    (80,90) % Alta Sicurezza
    (20,40) % Basse Latenze
    (50,70) % Bilanciata
};

% --- PUNTO OTTIMALE ---
\addplot[only marks, mark=*, mark size=5pt, red] coordinates {
    (32,75) % GIST Ottimale
};

% --- ETICHETTE ---
\node[above right] at (axis cs:80,90) {\footnotesize Alta Sicurezza};
\node[above right] at (axis cs:20,40) {\footnotesize Basse Latenze};
\node[above right] at (axis cs:50,70) {\footnotesize Bilanciata};
\node[above right, red] at (axis cs:32,75) {\footnotesize\bfseries GIST Ottimale};

\end{axis}
\end{tikzpicture}
\caption{Superficie di Pareto bidimensionale per l'ottimizzazione multi-obiettivo dei parametri di sistema.}
\label{fig:tradeoff}
\end{figure}

\end{document}